%% kapa: exam=1 round=25 year=2014 period=1 form=A subject=economics subject_ko=경제원론 seq=39 number=79
%% source: Q-Net 감정평가사 기출문제 (cst00309) artlSeq=5101207 file=extracted/1008/5101207/2014년 제25회 감정평가사 1차시험 기출문제/1교시(A형).hwp
%% conversion: kordoc 4.24.51 -> tools/gen1.py (auto); review=proofread-sonnet; flags=none
\begin{question}{79}
단기 필립스곡선은 $\pi_{ t } = \pi^{ e } -0.5(u_{ t } -u^{ n } )$이다. 중앙은행이 실업률을 $u^{ n }$수준으로 달성하기 위한 방법으로 옳은 것은? (단, $\pi_{ t }$는 $t$기의 물가상승률, $\pi^{ e }$는 예상 물가상승률, $u_{ t }$는 $t$기의 실업률, $u^{ n }$은 자연실업률이다.)
\begin{choices}
  \item 통화량 증가율을 높이다가 예고 없이 갑자기 낮춘다.
  \item 통화량 증가율을 낮추다가 예고 없이 갑자기 높인다.
  \item 통화량 증가율을 일정하게 유지한다고 공표한 다음 그대로 지킨다.
  \item 통화량 증가율을 일정하게 유지한다고 공표한 다음 더 높은 수준으로 바꾼다.
  \item 통화량 증가율을 일정하게 유지한다고 공표한 다음 더 낮은 수준으로 바꾼다.
\end{choices}
\end{question}
\answer{3}
